% FORMALIZACION_NUEVA de la incidencia firmada de la carta Paley preexistente.
% La construcción comienza en exc:aw; no identifica por decreto sus banderas
% con el cociente de todas las historias TPK.
\paragraph{Construcción del portador desde la incidencia de seis posiciones.}
\label{ii:ico:construccion-firmada}
La matriz \(A_W\) de \eqref{exc:aw}, que interviene previamente en el
transporte angular de la sección~\ref{sec:ii-enlace-angular-witt}, proporciona
un antecedente explícito de la realización icosaédrica. Sea
\(C=\operatorname{bal}(A_W)\), donde
\(\operatorname{bal}(0)=0\), \(\operatorname{bal}(1)=1\) y
\(\operatorname{bal}(2)=-1\). El orden de sus coordenadas es
\(\mathcal U=(\infty,0,1,2,3,4)\). La coordenatización entera satisface
\(C=C^{\mathsf T}\), \(C^2=5I_6\) y \(\operatorname{Tr}C=0\).
Definimos su recubrimiento firmado por
\begin{equation}
 \mathcal D_C=\mathcal U\times\{+1,-1\},\qquad
 (u,\sigma)\sim_C(v,\tau)
 \ \Longleftrightarrow\ u\ne v\ \text{ y }\ \sigma\tau=C_{uv}.
 \label{ii:ico:incidencia-firmada}
\end{equation}
La orientación duplica cada coordenada; la matriz decide las incidencias
entre las dos fibras correspondientes. La involución de este recubrimiento
es \(\jmath_C(u,\sigma)=(u,-\sigma)\).

\Needspace{12\baselineskip}
\begin{samepage}
\begin{proposition}[Realización isométrica de la incidencia firmada]
\label{ii:ico:realizacion-firmada}
El operador y los vectores
\begin{equation}
 E_C=\frac12\left(I_6+\frac C{\sqrt5}\right),\qquad
 \mathcal H_C=\operatorname{im}E_C,\qquad
 v_{u,\sigma}=\sigma\sqrt2 E_Ce_u
 \label{ii:ico:proyector-seis}
\end{equation}
realizan \(\mathcal D_C\) como doce vectores unitarios distintos en un
espacio euclidiano de dimensión tres. Existe una isometría explícita
\(T_C:\mathcal H_C\to\mathbb R^3\) para la que
\(\iota_C(u,\sigma)=T_Cv_{u,\sigma}\) es una biyección sobre
\(\Omega_{12}^{\rm ico}\), conserva la incidencia y satisface
\(\iota_C\jmath_C=-\iota_C\).
\end{proposition}
\end{samepage}
\begin{proof}
De \(C^2=5I_6\) se obtiene \(E_C^2=E_C=E_C^*\); su traza es tres.
Las identidades métricas son
\[
 \langle v_{u,\sigma},v_{v,\tau}\rangle
 =\begin{cases}\sigma\tau,&u=v,\\
 \sigma\tau C_{uv}/\sqrt5,&u\ne v.
 \end{cases}
\]
Así, los doce vectores son distintos y su incidencia de producto escalar
\(1/\sqrt5\) es exactamente~\eqref{ii:ico:incidencia-firmada}.
Para explicitar la isometría, sean las columnas de \(B\), en el orden
\(\mathcal U\),
\[
 B=\frac1\nu
 \begin{pmatrix}
 0&-1&-\varphi_g&0&\varphi_g&1\\
 -1&-\varphi_g&0&1&0&-\varphi_g\\
 -\varphi_g&0&-1&-\varphi_g&-1&0
 \end{pmatrix}.
\]
La sustitución de \(\varphi_g^2=\varphi_g+1\) da
\(B^*B=I_6+C/\sqrt5=2E_C\). Por tanto
\(T_C=B/\sqrt2\big|_{\mathcal H_C}\) es una isometría, y
\(T_Cv_{u,\sigma}=\sigma Be_u\). Las seis columnas de \(B\) y sus
opuestas son precisamente los doce vértices de
\eqref{ii:pent:eq:gravity-icosahedron-vertices}. La última identidad
demuestra también la compatibilidad de las involuciones.
\end{proof}

La realización dodecafásica precedente construye los doce vectores
directamente desde el proyector \(P_3^{\mathrm{ico}}\) y las posiciones del calendario.
El recubrimiento firmado proporciona otra expresión coordenada de ese
portador. Para el lector \(\pi_{\rm ico}\) del teorema de descenso,
la relación entre ambas expresiones se escribe
\begin{equation}
 X_{\rm TPK}^{\rm exc}\xrightarrow{\ r_\pm\ }
 \mathcal D_C\xrightarrow[\simeq]{\ \iota_C\ }
 \Omega_{12}^{\rm ico},\qquad
 \pi_{\rm ico}(x)=\sigma(x)Be_{u(x)},\quad
 r_\pm(x)=(u(x),\sigma(x)).
 \label{ii:ico:factorizacion-lector}
\end{equation}
Como \(\iota_C\) es una biyección, esta factorización determina
\(r_\pm=\iota_C^{-1}\pi_{\rm ico}\). Por tanto, \(r_\pm\) expresa
en la coordenatización firmada el lector que se esté utilizando: su introducción
no añade una hipótesis independiente para construir el portador
dodecafásico. La tabla de correspondencias precedente da esa expresión
para sus doce posiciones. Las propiedades (i)--(iv) conservan su
alcance propio cuando se identifica la realización con el cociente
de la ruta excepcional completa. En particular, el transporte conjunto
de marco, proyector e incidencia demuestra la covariancia entre coordenatizaciones;
la conservación de una incidencia fija por una transformación activa
es una identidad distinta. La orientación y la memoria permanecen
en el estado fuente durante ambas lecturas.
Una permutación firmada que transporta \(C\) transporta conjuntamente
\(E_C\), sus vectores y su incidencia; esta covariancia permite cotejar
las coordenatizaciones conservando la orientación de cada fibra.
